\section{Conclusion}\label{sec:conclusion}

Hierarchy equivalence is a partition of the vertices, and that partition is the right unit for indexing every clique-core hierarchy of a graph at once.
Storing chains instead of vertices, aligning labels to chains, ordering chains lexicographically and keeping each size's DFS array as runs with entry points gives an index that is \ratiomin$\times$ to \ratiomax$\times$ smaller than one tree per size, locates any community in tens of nanoseconds and lists it at memory-copy speed.
The construction is one all-size peel and one pass per size; its memory is the peel's, not the index's.
